% ============================================================================
%  chapter16.tex
%  Durable Storage: An Append-Only Log Database
% ============================================================================

\chapter{Durable Storage: An Append-Only Log Database}

\section{Why Persistence Is Hard}

\textbf{What persistence is.} A piece of data is \emph{persistent} when
it survives longer than the program that created it---across a program
restart, across a reboot, across a power loss. Data in RAM is
\emph{volatile}: it lives exactly as long as the machine's power and
the process's lifetime, and both can end at any instant. Data in a
file is \emph{persistent}: it lives on storage, independent of any
program and of the power state of the moment it was written. The whole
discipline of this chapter is about moving data across that line---RAM
to storage---without ever letting the transition lose the truth.

\begin{verbatim}
   Persistent vs volatile, the single idea:

     volatile   dies when the power dies          (RAM, your heap)
     persistent outlives the program and the      (files on disk)
                power cycle

   Your data is only as durable as its least persistent home:
   a "saved" file still being held in the page cache is, until
   fsync, one power loss away from vanishing (Section 16.3).
\end{verbatim}

Everything so far has stored data in \emph{RAM}, and RAM has a fatal
flaw for a database: it forgets. Turn the machine off, and the heap,
the hash tables, and the B-trees of earlier chapters are gone. A real
store must be \emph{durable}: the data must survive a power loss and
a crash, in a form that can be read back tomorrow.

\begin{verbatim}
   RAM vs storage, the two worlds:

     RAM          fast, volatile: fine while the machine runs
     file         slower, persistent: survives power loss

   The database's job is to shuttle between the two worlds:
   keep a fast copy in memory for answering queries NOW, and keep a
   safe copy in a file for answering queries AFTER the next boot.

   The hard part is the crash. A program can check every return
   value and still vanish at any instant. The database must be
   designed so that "vanished at any instant" leaves the file in a
   state from which the truth can be recovered.
\end{verbatim}

This chapter builds a real, durable key-value store: an
\emph{append-only log}. It is the smallest honest database---and the
same idea, under the names write-ahead log (WAL) and journal, sits
inside every production database you have ever used. When you finish
you will understand how data survives crashes, why you cannot write a
struct to a file and be done, and what ``durable'' actually means.

\begin{osconcept}
Durability is an operating system promise, and it is not free. When
your program calls \code{write}, the bytes do not reach the disk; they
reach the kernel's \emph{page cache}. Only \code{fsync} forces them
out. The kernel exists between your program and the hardware, and it
buffers for performance---so ``the file contains my data'' is a claim
the OS, not your program, ultimately guarantees. A durable database is
a careful conversation with the page cache.
\end{osconcept}

\section{The Append-Only Log}

The simplest durable structure is also the best one. Instead of
modifying a file in place, a database appends every change as a
\emph{record} to the end of a log. The log is the source of truth; the
in-memory structures are merely a cache of it. To rebuild after a
crash, you \emph{replay} the log: read every record and re-apply it.

\begin{verbatim}
   The append-only log, drawn:

   time --->  file "mydb.log"
              +--------+--------+--------+--------+
              | SET a=1 | SET b=2 | DEL a  | SET b=5|
              +--------+--------+--------+--------+
              records appended strictly in order, never overwritten

   Startup replay: read every record, apply it in order.
     SET a=1  ->  table[a] = 1
     SET b=2  ->  table[b] = 2
     DEL a    ->  remove a
     SET b=5  ->  table[b] = 5

   Result: { b = 5 }. The log reconstructs the exact final state.

   Why appending is the right shape:
     - appends are SEQUENTIAL writes, the fastest kind the disk does;
     - an append never destroys existing data, so a crash can only
       leave a PARTIAL RECORD at the tail -- never a corrupted middle;
     - replay is deterministic: same log, same state, every time.
\end{verbatim}

The property to notice is what a crash can and cannot do. Because we
only ever write at the tail, a crash mid-write can leave a torn record
at the end---but it cannot corrupt records that were already
completed. The database's recovery rule is therefore brutally simple:
\emph{replay from the start, and stop at the first damaged record.}
Everything before the damage is true.

\begin{verbatim}
   A crash at time X leaves:

        +--------+--------+--------+------XXXXX
        | SET a=1 | SET b=2 | SET a=1|  torn tail
        +--------+--------+--------+------XXXXX
                            ^ replay stops here: "a" and "b" are safe,
                              the torn record is discarded

   That is the entire crash story: replay the clean prefix, ignore the
   torn tail. No recovery magic, no repair pass -- the design makes
   recovery a read.
\end{verbatim}

\begin{osconcept}
Append-only is crash-friendly because the kernel's behavior matches
the structure: writes to a file are ordered, and once a write is
acknowledged and flushed, those bytes are durable. By never modifying
old bytes, the database ensures the past is always intact. Deciding
``the log is truth, memory is cache'' is a systems-design decision that
converts crash recovery from a ceremony into a read loop.
\end{osconcept}

\section{The Page Cache and fsync}

Before the record format, the durability mechanic deserves its own
section, because it decides what ``written'' means. When you call
\code{write(fd, buf, n)}, the kernel copies your bytes into the
\emph{page cache}---RAM that holds the file's pages---and returns. The
disk has not necessarily seen them. The kernel flushes cached pages
later, batched for efficiency. \code{fsync(fd)} is the syscall that
says \emph{now}: do not return until the file's data has actually
reached the storage device.

\begin{verbatim}
   Where a write() goes, layer by layer:

     your buffer  --write()-->  kernel page cache  --fsync()-->  disk
        ^                                ^
        "acknowledged"               "really there"
        (fast, not durable)          (slow, durable)

   The three-layer model:
     write()   copies bytes into the kernel's cache. Fast. Volatile:
               a power loss may drop it.
     fsync()   pushes the file's dirty pages to disk. Slow. Durable:
               after it returns, the data survives a power loss.

   Consequence for a database: a record is only durable AFTER fsync.
   A database that writes and never fsyncs is a database that lies
   to itself -- it has told the user "stored" while the truth still
   sits in volatile RAM.
\end{verbatim}

\begin{ruleofthumb}[What durable means: fsync before you claim success]
If you tell the user ``stored,'' the record must have been \code{fsync}'d.
For a group of changes that must survive together, append them all,
then \code{fsync} once. The ordering---write the record, then force it
out---is the entire discipline of crash-consistent storage.
\end{ruleofthumb}

\begin{osconcept}
The page cache is the kernel's buffer between your program and the
hardware, and it is exactly the Chapter~5 ``buffer'' idea at system
scale. The OS batches writes because disk I/O is millions of times
slower than RAM. \code{fsync} is the escape hatch that buys
durability when you need it. Deciding \emph{when} to pay that price is
the difference between a toy that loses data and a database that does
not.
\end{osconcept}

\section{The Record Format: Bytes, Not Structs}

A record is a sequence of bytes on disk. Here is the first subtlety:
you cannot simply \code{fwrite} a C struct. A struct contains
\emph{padding} (Chapter~3) whose size differs between machines and
compilers, its integers are stored in the machine's native
\emph{endianness} (Chapter~1), and reading it back on another machine
or with another compiler would misread every field. A durable format
is a \emph{wire format}: a precise, byte-by-byte description written
and read by hand.

\begin{verbatim}
   Our record, byte by byte:

   +---------+------+--------+--------+-----+-------+----------+
   | magic   | type | keylen | vallen | key | value | checksum |
   | 4 bytes | 1 B  | 1 B    | 4 B    | ... | ...   | 4 bytes  |
   +---------+------+--------+--------+-----+-------+----------+

   magic     "!DB1" = 0x31424421, detects a wrong file or a torn
            read at the record boundary
   type      1 = SET (store key = value), 2 = DEL (remove key)
   keylen    1 byte: key length 0..255
   vallen    4 bytes: value length
   key       keylen bytes
   value     vallen bytes
   checksum  a hash over everything before it, to detect corruption

   The total record is 10 + keylen + vallen + 4 bytes. Every field's
   position is fixed by the format, never by the struct layout.
\end{verbatim}

Two helpers define the wire format's integers. We choose
\emph{little-endian}: byte 0 is the least significant. On an x86 or
ARM machine that matches the hardware, and writing the bytes out
explicitly makes the format portable by construction.

\begin{lstlisting}[caption={Serialization helpers: ints as bytes, on purpose.}]
#include <stdio.h>
#include <stdlib.h>
#include <string.h>
#include <fcntl.h>
#include <unistd.h>
#include <errno.h>

#define MAGIC 0x31424421u    /* "!DB1" */

#define REC_SET 1
#define REC_DEL 2

/* Write a 32-bit unsigned integer as 4 little-endian bytes. */
static void put_u32(unsigned char *b, unsigned int v) {
    b[0] = (unsigned char) (v & 0xff);
    b[1] = (unsigned char) ((v >> 8) & 0xff);
    b[2] = (unsigned char) ((v >> 16) & 0xff);
    b[3] = (unsigned char) ((v >> 24) & 0xff);
}

/* Read a 32-bit unsigned integer from 4 little-endian bytes. */
static unsigned int get_u32(const unsigned char *b) {
    unsigned int v;

    v = (unsigned int) b[0];
    v |= ((unsigned int) b[1]) << 8;
    v |= ((unsigned int) b[2]) << 16;
    v |= ((unsigned int) b[3]) << 24;
    return v;
}

/* FNV-1a: a fast, simple 32-bit hash. We seed it so a checksum can
   be computed over several buffers as one rolling value.        */
static unsigned int fnv1a_seed(unsigned int h,
                               const unsigned char *p, size_t n) {
    size_t i;

    for (i = 0; i < n; i++) {
        h = h ^ p[i];
        h = h * 16777619u;
    }
    return h;
}

static unsigned int fnv1a(const unsigned char *p, size_t n) {
    return fnv1a_seed(2166136261u, p, n);
}
\end{lstlisting}

The checksum is the record's honesty check. We hash the header, the
key, and the value; a torn write, a flipped bit, or a truncated tail
changes the hash, and replay detects it. A checksum does not repair
data---it makes \emph{corruption detectable}, which is what lets
replay stop safely at the first bad record instead of trusting it.

\begin{osconcept}
A wire format is the boundary between two representations of the same
data: in-memory structs (padding, native endianness, pointers) and
durable bytes (no padding, fixed byte order, no pointers). Only the
kernel delivers your bytes to the disk, but \emph{you} define what
they mean. Getting this boundary explicit is the same discipline we
will need for the network in Chapter~17---a network protocol is a wire
format that travels over a socket instead of a file.
\end{osconcept}

\section{The Write Path: Append and fsync}

The write path has two responsibilities: put a well-formed record at
the end of the log, and make it durable before reporting success. The
first needs a helper that writes \emph{all} requested bytes (a plain
\code{write} may write fewer), and the second needs \code{fsync}.

\begin{lstlisting}[caption={write\_all: loop until every byte is written.}]
/* Write all n bytes, retrying on EINTR. Returns 0 or -1. */
static int write_all(int fd, const void *buf, size_t n) {
    const unsigned char *p;
    size_t written;

    p = buf;
    written = 0;
    while (written < n) {
        ssize_t r;

        r = write(fd, p + written, n - written);
        if (r < 0) {
            if (errno == EINTR) {
                continue;        /* a signal interrupted us; retry */
            }
            return -1;
        }
        written = written + (size_t) r;
    }
    return 0;
}
\end{lstlisting}

The append function builds the record into its on-disk form (header,
key, value, checksum), writes the pieces in order, and fsyncs. The
order of the bytes on disk is fixed by the format; the order of the
fsync is fixed by the discipline: \emph{write first, claim success
only after fsync}.

\begin{lstlisting}[caption={append\_record: one durable record at the tail.}]
typedef struct dblog dblog;   /* defined in the next section */

/* Append one record. Returns 0 on success, -1 on failure. */
static int append_record(dblog *db, int type,
                         const char *key, int keylen,
                         const char *value, int vallen) {
    unsigned char header[10];
    unsigned char check[4];
    unsigned int sum;

    put_u32(header, MAGIC);
    header[4] = (unsigned char) type;
    header[5] = (unsigned char) keylen;
    put_u32(header + 6, (unsigned int) vallen);

    /* Checksum over header + key + value, rolled into one hash. */
    sum = fnv1a(header, 10);
    sum = fnv1a_seed(sum, (const unsigned char *) key, (size_t) keylen);
    sum = fnv1a_seed(sum, (const unsigned char *) value, (size_t) vallen);
    put_u32(check, sum);

    /* The four pieces of the record, in the wire order. */
    if (write_all(db->fd, header, 10) < 0) {
        return -1;
    }
    if (write_all(db->fd, key, (size_t) keylen) < 0) {
        return -1;
    }
    if (write_all(db->fd, value, (size_t) vallen) < 0) {
        return -1;
    }
    if (write_all(db->fd, check, 4) < 0) {
        return -1;
    }

    /* Durable now: do not tell the caller "stored" until it is. */
    if (fsync(db->fd) < 0) {
        return -1;
    }

    db->offset = db->offset + 10 + keylen + vallen + 4;
    return 0;
}
\end{lstlisting}

\begin{verbatim}
   The order is the contract:

     1. append the record's bytes  (fast, into the page cache)
     2. fsync()                    (slow, durable -- the record is
                                    now safe against a power loss)
     3. only then update memory    (the in-memory table is a cache;
                                    the log is the truth)

   If the process crashes between 1 and 2, replay may find a torn
   tail and stop -- but it will never invent a record that was not
   fully written. The fsync is the line between "maybe" and "stored."
\end{verbatim}

\section{The Read Path: Replay}

Reading a log back is the mirror of appending, with one difference:
appending assumes the file is well-formed, while replay must be
\emph{suspicious}. A record may be torn, truncated, or corrupted, and
replay must detect it and stop.

The reader uses \code{pread}, a system call we have not needed before:
\code{pread(fd, buf, n, offset)} reads at a given offset \emph{without
moving the file position}. That is exactly right for replay---each
record is read at a position we track ourselves, and the shared file
position stays untouched.

\begin{lstlisting}[caption={read\_record: read one record, or say why you cannot.}]
/* Read one record at *pos. Returns:
     1  -> a valid record (key/value returned, *pos advanced)
     0  -> clean end of file
    -1  -> torn or corrupted (replay must stop here)             */
static int read_record(int fd, long *pos,
                       int *type, char **key, int *keylen,
                       char **value, int *vallen) {
    unsigned char header[10];
    unsigned char check[4];
    unsigned char *key_buf;
    unsigned char *value_buf;
    unsigned int want;
    unsigned int got;
    int n;

    n = (int) pread(fd, header, 10, *pos);
    if (n == 0) {
        return 0;                       /* clean end of file */
    }
    if (n != 10) {
        return -1;                      /* torn header */
    }
    if (get_u32(header) != MAGIC) {
        return -1;                      /* wrong file or bad boundary */
    }

    *type = header[4];
    *keylen = (int) header[5];
    *vallen = (int) get_u32(header + 6);

    key_buf = malloc((size_t) *keylen + 1);
    value_buf = malloc((size_t) *vallen);
    if (key_buf == NULL || value_buf == NULL) {
        free(key_buf);
        free(value_buf);
        return -1;
    }

    n = (int) pread(fd, key_buf, (size_t) *keylen, *pos + 10);
    if (n != *keylen) {
        free(key_buf);
        free(value_buf);
        return -1;
    }
    n = (int) pread(fd, value_buf, (size_t) *vallen,
                    *pos + 10 + *keylen);
    if (n != *vallen) {
        free(key_buf);
        free(value_buf);
        return -1;
    }

    n = (int) pread(fd, check, 4, *pos + 10 + *keylen + *vallen);
    if (n != 4) {
        free(key_buf);
        free(value_buf);
        return -1;
    }

    want = fnv1a(header, 10);
    want = fnv1a_seed(want, key_buf, (size_t) *keylen);
    want = fnv1a_seed(want, value_buf, (size_t) *vallen);
    got = get_u32(check);
    if (want != got) {
        free(key_buf);
        free(value_buf);
        return -1;                      /* checksum mismatch */
    }

    key_buf[*keylen] = '\0';
    *key = (char *) key_buf;
    *value = (char *) value_buf;
    *pos = *pos + 10 + *keylen + *vallen + 4;
    return 1;
}
\end{lstlisting}

Replay is then a single loop: read a record, apply it, continue; stop
at the first \code{0} (clean end) or \code{-1} (damage). The partially
rebuilt table is the truth; the records after the damage, if any, are
discarded.

\begin{lstlisting}[caption={db\_replay: read the clean prefix, apply it.}]
/* Rebuild the in-memory index by replaying the log from the start. */
static void db_replay(dblog *db) {
    long pos;
    int r;

    pos = 0;
    while (1) {
        int type;
        int keylen;
        int vallen;
        char *key;
        char *value;

        r = read_record(db->fd, &pos, &type, &key, &keylen,
                        &value, &vallen);
        if (r <= 0) {
            break;           /* clean end, or damage: stop here */
        }

        if (type == REC_SET) {
            table_put(db, key, keylen, value, vallen);
        } else if (type == REC_DEL) {
            table_del(db, key, keylen);
        }

        free(key);
        free(value);
    }
}
\end{lstlisting}

\begin{verbatim}
   Replay is a loop over a file, and that is the whole point:
   the recovery procedure is the same code as startup. There is no
   "repair mode," no interactive recovery, no special case -- the
   log was designed so that reading it IS recovering it.
\end{verbatim}

\section{The In-Memory Index and the API}

The in-memory half is the hash table of Chapter~4, with one important
role: it is a \emph{cache}, not the authority. Writes go to the log
first and then to the table; reads come from the table. A lookup never
touches the disk, which is why the database is fast; the log never
lies, which is why it is durable.

\begin{lstlisting}[caption={The database: a log plus a hash-table cache.}]
typedef struct kv {
    char *key;
    int keylen;
    char *value;
    int vallen;
    struct kv *next;
} kv;

typedef struct dblog {
    int fd;
    kv **buckets;      /* the hash table (Chapter 4) */
    int nbuckets;
    long count;        /* number of live entries      */
    long offset;       /* current end of the log      */
} dblog;

static int hash_key(const char *key, int keylen, int nbuckets) {
    return (int) (fnv1a((const unsigned char *) key,
                        (size_t) keylen) % (unsigned int) nbuckets);
}

/* Look up a key in the table; returns the node or NULL. */
static kv *table_get(const dblog *db, const char *key, int keylen) {
    kv *cur;

    cur = db->buckets[hash_key(key, keylen, db->nbuckets)];
    while (cur != NULL) {
        if (cur->keylen == keylen &&
            memcmp(cur->key, key, (size_t) keylen) == 0) {
            return cur;
        }
        cur = cur->next;
    }
    return NULL;
}

/* Insert or update. Returns 1 if a new key was added. */
static int table_put(dblog *db, const char *key, int keylen,
                     const char *value, int vallen) {
    kv *node;
    kv *cur;
    int b;

    cur = table_get(db, key, keylen);
    if (cur != NULL) {
        /* Update the existing node in place. */
        char *new_value;

        new_value = malloc((size_t) vallen);
        if (new_value == NULL) {
            return 0;
        }
        memcpy(new_value, value, (size_t) vallen);
        free(cur->value);
        cur->value = new_value;
        cur->vallen = vallen;
        return 0;
    }

    node = malloc(sizeof(kv));
    if (node == NULL) {
        return 0;
    }
    node->key = malloc((size_t) keylen);
    node->value = malloc((size_t) vallen);
    if (node->key == NULL || node->value == NULL) {
        free(node->key);
        free(node->value);
        free(node);
        return 0;
    }
    memcpy(node->key, key, (size_t) keylen);
    memcpy(node->value, value, (size_t) vallen);
    node->keylen = keylen;
    node->vallen = vallen;

    b = hash_key(key, keylen, db->nbuckets);
    node->next = db->buckets[b];
    db->buckets[b] = node;
    db->count = db->count + 1;
    return 1;
}

/* Remove a key if present. Returns 1 if it was removed. */
static int table_del(dblog *db, const char *key, int keylen) {
    kv **cur;
    kv *victim;

    cur = &db->buckets[hash_key(key, keylen, db->nbuckets)];
    while (*cur != NULL) {
        if ((*cur)->keylen == keylen &&
            memcmp((*cur)->key, key, (size_t) keylen) == 0) {
            victim = *cur;
            *cur = victim->next;
            free(victim->key);
            free(victim->value);
            free(victim);
            db->count = db->count - 1;
            return 1;
        }
        cur = &(*cur)->next;
    }
    return 0;
}
\end{lstlisting}

The public API is three functions, each a short story about the log
and the cache:

\begin{lstlisting}[caption={The API: log first, then cache.}]
/* Store key = value durably. */
int db_set(dblog *db, const char *key, int keylen,
           const char *value, int vallen) {
    if (append_record(db, REC_SET, key, keylen, value, vallen) < 0) {
        return -1;          /* the log did not take the record */
    }
    table_put(db, key, keylen, value, vallen);
    return 0;
}

/* Remove key durably. Returns 1 if it existed, 0 if not. */
int db_del(dblog *db, const char *key, int keylen) {
    int existed;

    existed = (table_get(db, key, keylen) != NULL);
    if (existed) {
        if (append_record(db, REC_DEL, key, keylen, NULL, 0) < 0) {
            return -1;
        }
        table_del(db, key, keylen);
    }
    return existed;
}

/* Fetch a key. On success returns 1 and a malloc'd copy of the
   value (the caller frees it); returns 0 if absent.             */
int db_get(const dblog *db, const char *key, int keylen,
           char **out_value, int *out_vallen) {
    kv *node;
    char *copy;

    node = table_get(db, key, keylen);
    if (node == NULL) {
        return 0;
    }
    copy = malloc((size_t) node->vallen);
    if (copy == NULL) {
        return 0;
    }
    memcpy(copy, node->value, (size_t) node->vallen);
    *out_value = copy;
    *out_vallen = node->vallen;
    return 1;
}
\end{lstlisting}

Notice the ordering in \code{db\_set}: the record goes to the log
\emph{first}, and the table is updated only if the append succeeded.
If the append fails, memory must not pretend the write happened. If
the append succeeds and we crash before touching the table, replay
rebuilds it. \emph{Log first, cache second} is the rule, and it is
what makes the two halves consistent.

\begin{osconcept}
The log/cache split is a layering decision of the same shape as
Chapter~14's parser: an authority (the log) and a fast view (the
table), with one-directional trust. The kernel provides the durability
primitive (\code{fsync}); the database chooses when to pay for it; the
cache makes reads cheap. Every production database---SQLite's journal,
PostgreSQL's WAL---is this same split, tuned and enlarged.
\end{osconcept}

\section{Compaction: Shrinking the Log}

An append-only log only grows. Every update of a key appends another
record, so after a long life the log is full of dead records---old
versions of keys that were later changed or deleted. The cure is
\emph{compaction}: rewrite a fresh log containing only the \emph{live}
entries, then swap it in atomically. The compaction walks the in-memory
table (which already holds the truth) and writes one \code{SET} record
per live entry.

The swap uses \code{rename}, and that one call carries the whole
atomicity: \code{rename} either succeeds completely or changes nothing,
so at any instant a reader sees either the old log or the new one,
never a mixture.

\begin{lstlisting}[caption={db\_compact: live entries only, then an atomic swap.}]
/* Rewrite the log with only the live entries. Returns 0 or -1. */
int db_compact(dblog *db, const char *path) {
    char tmp[1024];
    int new_fd;
    int b;
    kv *cur;

    snprintf(tmp, sizeof(tmp), "%s.tmp", path);
    new_fd = open(tmp, O_WRONLY | O_CREAT | O_TRUNC, 0644);
    if (new_fd < 0) {
        return -1;
    }

    /* Walk every bucket; write one SET record per live entry. */
    for (b = 0; b < db->nbuckets; b++) {
        cur = db->buckets[b];
        while (cur != NULL) {
            unsigned char header[10];
            unsigned char check[4];
            unsigned int sum;

            put_u32(header, MAGIC);
            header[4] = REC_SET;
            header[5] = (unsigned char) cur->keylen;
            put_u32(header + 6, (unsigned int) cur->vallen);

            sum = fnv1a(header, 10);
            sum = fnv1a_seed(sum, (const unsigned char *) cur->key,
                             (size_t) cur->keylen);
            sum = fnv1a_seed(sum, (const unsigned char *) cur->value,
                             (size_t) cur->vallen);
            put_u32(check, sum);

            if (write_all(new_fd, header, 10) < 0) {
                goto fail;
            }
            if (write_all(new_fd, cur->key,
                          (size_t) cur->keylen) < 0) {
                goto fail;
            }
            if (write_all(new_fd, cur->value,
                          (size_t) cur->vallen) < 0) {
                goto fail;
            }
            if (write_all(new_fd, check, 4) < 0) {
                goto fail;
            }

            cur = cur->next;
        }
    }

    /* Durable, then atomic: the rename is all-or-nothing. */
    if (fsync(new_fd) < 0) {
        goto fail;
    }
    close(new_fd);

    if (rename(tmp, path) < 0) {
        return -1;
    }

    close(db->fd);
    db->fd = open(path, O_RDWR, 0644);
    if (db->fd < 0) {
        return -1;
    }
    db->offset = lseek(db->fd, 0, SEEK_END);
    if (db->offset < 0) {
        return -1;
    }
    return 0;

fail:
    close(new_fd);
    unlink(tmp);
    return -1;
}
\end{lstlisting}

The \code{goto fail} deserves an explanation, because it is the one
place this book uses \code{goto} on purpose. In C, error handling has
no \code{finally} clause; the idiomatic pattern is a single cleanup
label that every error path jumps to, so that cleanup code exists
\emph{once} instead of being repeated after every failing call.
Every \code{goto fail} means ``something went wrong; clean up the
temporary file and return failure.''

\begin{verbatim}
   Compaction in pictures:

     before:  log = SET a=1, SET b=2, SET a=1, SET a=3, DEL b
              (5 records; only {a=3} is live)
     compact: walk the table, write SET a=3 to "db.tmp", fsync
     swap:    rename("db.tmp", "db")  -- atomic
     after:   log = SET a=3  (one record; the dead history is gone)

   A crash at any point leaves either the old (bigger) log or the new
   (smaller) one. Both are valid; the rename is the boundary between
   them, and it is atomic.
\end{verbatim}

\begin{osconcept}
Atomic rename is a kernel guarantee that the database leans on: the
filesystem makes the name switch indivisible, so compaction never
exposes a half-written file. This is the same ``reference counting and
naming is kernel business'' lesson as Chapter~8, applied to the most
delicate moment in a database's life. The kernel's atomicity is the
only reason swap-in-place works at all.
\end{osconcept}

\section{Practical: The Whole Database}

Let us assemble every piece---serialization, append, replay, index,
compaction---into one program, plus the two glue functions we have not
shown (\code{db\_open} opens or creates the log and replays it;
\code{db\_close} frees the index). The pieces have all appeared above;
here is how they fit.

\begin{lstlisting}[caption={dblog.c: opening, closing, and a demo main.}]
int db_open(dblog *db, const char *path) {
    int fd;

    fd = open(path, O_RDWR | O_CREAT, 0644);
    if (fd < 0) {
        return -1;
    }
    db->fd = fd;
    db->nbuckets = 1024;
    db->count = 0;
    db->offset = lseek(fd, 0, SEEK_END);
    if (db->offset < 0) {
        close(fd);
        return -1;
    }
    db->buckets = calloc((size_t) db->nbuckets, sizeof(kv *));
    if (db->buckets == NULL) {
        close(fd);
        return -1;
    }

    db_replay(db);        /* rebuild the cache from the log */
    return 0;
}

void db_close(dblog *db) {
    int b;
    kv *cur;

    for (b = 0; b < db->nbuckets; b++) {
        cur = db->buckets[b];
        while (cur != NULL) {
            kv *next;

            next = cur->next;
            free(cur->key);
            free(cur->value);
            free(cur);
            cur = next;
        }
    }
    free(db->buckets);
    close(db->fd);
    db->buckets = NULL;
    db->fd = -1;
}

int main(int argc, char *argv[]) {
    dblog db;
    char *value;
    int vallen;

    if (argc != 2) {
        fprintf(stderr, "usage: %s dbfile\n", argv[0]);
        return 1;
    }

    if (db_open(&db, argv[1]) < 0) {
        perror("db_open");
        return 1;
    }

    db_set(&db, "answer", 6, "42", 2);
    db_set(&db, "name", 4, "ada", 3);
    db_set(&db, "answer", 6, "43", 2);     /* update: appends a record */
    db_del(&db, "old", 3);                 /* absent: no record written */

    if (db_get(&db, "answer", 6, &value, &vallen) == 1) {
        printf("answer = %.*s\n", vallen, value);
        free(value);
    }

    db_compact(&db, argv[1]);
    db_close(&db);
    return 0;
}
\end{lstlisting}

\begin{verbatim}
   Run it, then look at the file -- the wire format is right there:

     $ ./dblog /tmp/mydb.log
     answer = 43
     $ xxd /tmp/mydb.log
     00000000: 21 44 42 31 01 06 02 00 00 00   !DB1.........
                ^^ ^^ ^^ ^^ ^^    ^^^^^^^^
                magic  type keylen vallen (2, little-endian)
                "!DB1"  SET  = 6
     0000000a: 61 6e 73 77 65 72               "answer"
     00000010: 34 32                           "42"
     00000012: XX XX XX XX                      checksum
     ... and so on for each record.

   Re-run the program: it opens the same file, REPLAYS it, and
   prints answer = 43 again -- the data survived the previous run.
   That is durability, demonstrated by reading the file back.
\end{verbatim}

\begin{verbatim}
   To watch recovery from a torn tail, append garbage:

     $ echo "garbage" >> /tmp/mydb.log
     $ ./dblog /tmp/mydb.log        (the magic check fails; replay
     answer = 43                     stops before the garbage, and
                                     the clean prefix is intact)

   No crash simulator needed: the checksum and magic find the damage,
   replay stops, and the truth before the damage is preserved.
\end{verbatim}

\section{Where the Real Databases Go From Here}

This log is a real, durable, replayable store in about three hundred
lines. Production databases take the same skeleton and add four things,
each of which we have met in earlier chapters:

\begin{verbatim}
   what real systems add            where we met the idea
   ----------------------------     -------------------------------
   batched, asynchronous fsync      the page cache (this chapter)
   (group commit)                   and Chapter 5's buffering
   sorted on-disk storage           the B-tree family (Chapter 4)
   instead of replay-from-zero
   concurrency and locking          threads (Chapter 11)
   a network front end              the next chapter
\end{verbatim}

The lesson is that durability is not a library---it is a discipline
of \emph{fsync ordering}, \emph{checksums}, and \emph{atomic
renames}, exactly the ones you just wrote. Everything else is
performance and convenience built on top.

\begin{researchnote}[Log-structured storage, from the source]
The decision to make a log the heart of a storage system has a famous
pedigree. In 1991, Mendel Rosenblum and John Ousterhout published
``The Design and Implementation of a Log-Structured File System,'' the
paper behind the LFS filesystem---and, later, much of the modern
flash-storage world. Its central idea is precisely this chapter's:
\emph{treat the storage medium as one continuous log, write
sequentially, and use a compacting cleaner to reclaim dead space}.

\begin{verbatim}
   The paper's thesis, condensed:

       "By making all writes sequential appends to a log, the disk's
        expensive seeks largely disappear, and a cleaner process that
        rewrites only live data keeps the log from growing without
        bound."

   Translate to this chapter: append (sequential writes), checksum
   (detect damage), compact (the cleaner). Your dblog is a
   log-structured store in miniature, and the LFS paper is the
   original argument for why that shape is fast and crash-safe.
\end{verbatim}

Reading it is worthwhile for one reason above all: it shows that the
``simple'' log you just built was, in 1991, a research breakthrough,
and that its ideas now sit inside everything from SSDs to databases.

\paperref{M.~Rosenblum, J.~K.~Ousterhout, ``The Design and
Implementation of a Log-Structured File System,'' ACM SOSP 1991}{https://web.stanford.edu/~ouster/cgi-bin/papers/lfs.pdf}
\end{researchnote}

\section{Structures in Bytes}

The chapter's whole byte structure is the record format, so this is
the page to keep. One record, with absolute offsets:

\begin{verbatim}
   offset  size  field           example (SET answer = "42")
   0       4     magic "!DB1"    21 44 42 31
   4       1     type            01 = SET
   5       1     keylen          06
   6       4     vallen          02 00 00 00   (little-endian 2)
   10      keylen key            "answer"
   10+keylen vallen  value       "42"
   10+keylen+vallen  4  checksum  XX XX XX XX  (FNV-1a)

   total length = 10 + keylen + vallen + 4.
   the file is just records back to back:

   [SET a=1][SET b=2][DEL a][SET b=5]    append-only, never rewritten
\end{verbatim}

Little-endian is a deliberate choice, written byte by byte by
\code{put\_u32} and read back by \code{get\_u32}; the checksum covers
header + key + value, so a torn tail fails the hash and replay stops
at the last good byte. Compaction rewrites the same records into a
temp file and \code{rename}s it over the original---the on-disk
structure is unchanged, only shorter.

\section{Summary}

\begin{verbatim}
   The durable-storage toolbox, complete:

   the page cache      where write() actually puts your bytes
   fsync()             the line between "maybe" and "stored"
   append-only log     sequential writes; the past is never changed
   wire format         bytes, not structs: fixed order, explicit
                       endianness, no padding
   checksum            makes corruption detectable
   replay              reading the log IS recovering the database
   pread()             read at an offset without moving the position
   compaction          rewrite only live entries
   rename()            the atomic swap between old log and new

   Rules that keep durable code correct:
     1. write the record, THEN fsync, THEN claim success;
     2. append only at the tail; never overwrite old bytes;
     3. serialize with a fixed wire format -- never fwrite a struct;
     4. on any damage, replay stops: the clean prefix is the truth;
     5. compact to a temp file, fsync, and rename atomically.
\end{verbatim}

\begin{osconcept}
Durability is the operating system's oldest promise---\emph{your data
will still be there tomorrow}---and this chapter was the conversation
that keeps it: asking the kernel's page cache to flush, trusting its
atomic renames, and using its files as the wire format's transport.
The log is where the OS's storage services meet the data structures of
the first half of this book, and it is the same architecture every
real database runs.
\end{osconcept}

The database is durable on disk, and it has a command language. The
final chapter gives it a voice across the machine: we build sockets,
learn the TCP and UDP protocols, and turn this store---and the event
loop of Chapter~10---into a server that answers the world.


